\theory{Field theory}{What is a number system in which the four basic arithmetic operations work, and how can one enlarge such a system without losing control of equations?}{Field theory studies fields and their extensions. It explains why some polynomial equations can be solved by familiar numbers, why others require new numbers, and how the symmetries of the solutions reveal the answer.}{Algebraic equations, finite-field arithmetic, coding, cryptography, algebraic geometry, number theory, and mathematical physics.}{Field extensions, minimal polynomials, and splitting fields determine which algebraic operations are possible over a chosen number system.}

\paragraph{Évariste Galois}
Galois made field extensions and their automorphism groups central to the study of polynomial equations. His correspondence between solvability by radicals and solvability of the associated Galois group explains why some equations admit formulas and others do not.

\paragraph{Richard Dedekind and finite-field researchers}
Dedekind clarified the arithmetic of field extensions through embeddings, norms, and algebraic integers. The systematic study of finite fields and extension degrees later supplied the exact arithmetic used in coding, cryptography, and algebraic number theory \cite{field_theory_ref1}.
\end{historylist}
\section*{3. Theory}
\subsection*{Core theory}
\paragraph{What the name means}
The phrase "field theory" has several meanings. In mathematics it means the study of fields, such as the rational numbers, real numbers, complex numbers, and finite fields. In physics, field theory studies quantities spread through space and time: classical field theory, quantum field theory, and statistical field theory are related physical subjects. A grand unified theory is a proposed physical framework joining interactions. In psychology, Kurt Lewin's field theory studies the interaction of a person and environment, while sociological field theory studies social actors and local social orders. This chapter develops the mathematical meaning while keeping these other uses distinct \cite{field_theory_ref1,field_theory_ref4}.

\paragraph{The basic number system}
A field is a set of numbers in which addition, subtraction, multiplication, and division by a nonzero number are always possible, and the familiar rules of arithmetic hold. The rational numbers $\mathbb{Q}$ form a field. The real numbers $\mathbb{R}$ and complex numbers $\mathbb{C}$ form larger fields. The integers do not form a field because, for example, $1/2$ is not an integer.

Finite fields also exist. The integers modulo a prime $p$, written $\mathbb{F}_p$, form a field because every nonzero residue has a multiplicative inverse. Modulo $5$, the inverse of $2$ is $3$, since $2\cdot3=6$ has remainder $1$. Arithmetic in finite fields is exact and bounded, which is why it is useful in digital communication and cryptography.

\paragraph{Extensions: making room for new solutions}
\bookfigure{field_galois_lattice}{A field-extension lattice records how larger number systems contain one another.}
If a field $K$ is contained in a larger field $L$, then $L/K$ is a field extension. The smallest field containing $K$ and an element $\alpha$ is denoted $K(\alpha)$. If every element of $L$ can be written as a combination of finitely many powers of $\alpha$ with coefficients in $K$, the extension is finite. Its degree $[L:K]$ is the number of independent coefficients required.

The equation $x^2-2=0$ has no rational solution. Adjoining $\sqrt2$ produces $\mathbb{Q}(\sqrt2)$, whose elements have the form $a+b\sqrt2$ with $a,b\in\mathbb{Q}$. Since $(\sqrt2)^2=2$, every higher power reduces to one of the first two powers, so the degree is $2$. This is the controlled way in which algebra enlarges the number system.

An element $\alpha$ is algebraic over $K$ if it satisfies a nonzero polynomial with coefficients in $K$. Among all monic polynomials that vanish at $\alpha$, one of least degree is the minimal polynomial. It is irreducible over $K$ and records all algebraic relations that $\alpha$ has over the base field. If no such polynomial exists, $\alpha$ is transcendental. The number $\pi$ is an important example of a transcendental number over $\mathbb{Q}$.

\paragraph{Splitting fields and symmetries}
Given a polynomial $f(x)$ over $K$, its splitting field is the smallest extension in which $f$ factors completely into linear factors. For $x^2-2$, the splitting field is $\mathbb{Q}(\sqrt2)$. For $x^3-2$, the field must contain the real cube root of $2$ and the complex cube roots, so it also contains a primitive cube root of unity.

The automorphisms of a splitting field that fix every element of $K$ form its Galois group. An automorphism is a reversible map preserving addition and multiplication. It may permute the roots of the polynomial, but it must preserve every algebraic relation among them. For $x^2-2$, the two automorphisms send $\sqrt2$ to $\sqrt2$ or $-\sqrt2$. The group therefore has two elements.

Galois theory creates a dictionary: intermediate fields between $K$ and a suitable extension correspond to subgroups of the Galois group. Larger subgroups fix smaller fields, and smaller subgroups fix larger fields. The dictionary changes the problem of solving equations into the problem of understanding symmetries. It explains why general equations of degree five cannot be solved by radicals, whereas equations of degrees two, three, and four can.

\paragraph{Separable and normal extensions}
An algebraic extension is separable when every element has a minimal polynomial with distinct roots. In characteristic zero, every algebraic extension is separable. In positive characteristic, repeated roots can occur. For example, over a field of characteristic $p$, the polynomial $x^p-a$ has derivative zero and may have only one repeated root in a suitable extension.

An extension is normal when every irreducible polynomial over the base field that has one root in the extension splits completely there. A finite extension that is both normal and separable is Galois. The distinction matters because the clean subgroup dictionary belongs to Galois extensions, while inseparable extensions behave differently.

\paragraph{Constructing fields and understanding their size}
Every field has a characteristic. If adding $1$ repeatedly eventually gives $0$, the characteristic is a prime $p$ and the field contains a copy of $\mathbb{F}_p$. Otherwise its characteristic is zero and it contains a copy of $\mathbb{Q}$. A finite field has exactly $p^n$ elements for a prime $p$ and positive integer $n$, and any two fields with the same number of elements are isomorphic. The multiplicative nonzero elements of a finite field form a cyclic group \cite{field_theory_ref2,field_theory_ref3}.

Field extensions can be built by quotienting a polynomial ring by an irreducible polynomial. For example,
\[
\mathbb{Q}[x]/(x^2-2)
\]
behaves like $\mathbb{Q}(\sqrt2)$: the class of $x$ acts as $\sqrt2$, and the relation $x^2=2$ is imposed. Finite fields are constructed similarly from $\mathbb{F}_p[x]$ using irreducible polynomials.

\paragraph{Geometry, number theory, and computation}
Field theory is the algebraic foundation of algebraic geometry. Coordinates of points on a variety form a field or a ring, and extending the field describes new points that become visible after allowing new coordinates. In number theory, number fields are finite extensions of $\mathbb{Q}$; their rings of integers, ideals, places, and Galois groups encode arithmetic information.

Finite fields make arithmetic possible in computers without rounding error. Reed--Solomon codes evaluate polynomials over a finite field and recover the original message when some symbols are corrupted. Elliptic-curve cryptography uses the group of points on an elliptic curve over a finite field; the security comes from the difficulty of reversing a repeated group operation. A field is not itself a cryptosystem or code, but supplies the exact arithmetic on which these systems rely.

\paragraph{What the theory depends on and where it is useful}
Field theory depends on set-based algebra, polynomial arithmetic, vector spaces, and group theory. It supplies the coefficient systems used by ring theory, representation theory, algebraic geometry, number theory, coding theory, and parts of topology. Galois groups connect field theory to group theory; splitting fields connect it to polynomial equations; finite fields connect it to algorithms.

The theory also has limits. A field extension can tell us that roots exist and how they are related without producing a convenient decimal formula. A polynomial may be solvable by radicals in principle while its expression is too large to use. In positive characteristic, separability and derivative arguments require extra care.

\section*{4. Important theorems and results}
\begin{enumerate}[leftmargin=*,itemsep=0.35em]
\item \textbf{Tower law.} If $K\subseteq L\subseteq M$ are finite extensions, then
\[
[M:K]=[M:L][L:K].
\]
It is the arithmetic rule for counting independent coefficients.

\item \textbf{Primitive element theorem.} Many finite separable extensions can be generated by one element: $L=K(\alpha)$ \cite{field_theory_ref1}.

\item \textbf{Galois correspondence.} For a finite Galois extension, intermediate fields and subgroups of the Galois group correspond in reverse order \cite{field_theory_ref2}.

\item \textbf{Fundamental theorem of algebra.} Every nonconstant complex polynomial factors into linear factors over $\mathbb{C}$; thus $\mathbb{C}$ is algebraically closed.

\item \textbf{Finite-field theorem.} For every prime power $p^n$ there is, up to isomorphism, exactly one field with $p^n$ elements, and its multiplicative group is cyclic \cite{field_theory_ref3}.
\end{enumerate}

\theoryapplications
\theoryconnections{field theory}{%
\item Ring theory supplies polynomial quotients and ideals, linear algebra supplies extension degree, and Galois theory studies the symmetries of normal separable extensions.
\item Algebraic geometry and number theory provide geometric and arithmetic examples of fields.
}{%
\item Field theory helps construct solvable equations, finite fields for coding and cryptography, and coordinate fields in algebraic geometry.
\item It tells us which operations preserve a problem and when adjoining a new element creates genuinely new solutions.
}
\section*{Glossary}
\begin{description}[leftmargin=*,style=nextline,itemsep=0.3em]
\item[\textbf{Field extension.}] A larger field $L$ containing a smaller field $K$, written $L/K$, formed by adjoining new elements to $K$ and closing under arithmetic operations.
\item[\textbf{Minimal polynomial.}] The monic polynomial of lowest degree with coefficients in $K$ that vanishes at a given algebraic element; it encodes all algebraic relations that element has over $K$.
\item[\textbf{Splitting field.}] The smallest field extension over which a given polynomial factors completely into linear factors.
\item[\textbf{Galois group.}] The group of field automorphisms of a splitting field that fix every element of the base field; it records the symmetries among the roots.
\item[\textbf{Separable extension.}] An algebraic extension in which every element has a minimal polynomial with distinct roots; in characteristic zero every extension is separable.
\item[\textbf{Normal extension.}] An algebraic extension in which every irreducible polynomial over the base field that has one root in the extension splits completely there.
\item[\textbf{Algebraic element.}] An element satisfying a nonzero polynomial equation with coefficients in the base field, as opposed to a transcendental element like $\pi$.
\item[\textbf{Finite field.}] A field with finitely many elements, necessarily $p^n$ for a prime $p$; its nonzero elements form a cyclic multiplicative group.
\item[\textbf{Degree of an extension.}] The dimension $[L\!:\!K]$ of $L$ as a vector space over $K$.
\item[\textbf{Characteristic.}] The smallest positive integer $n$ such that $n\cdot1=0$ in a field, or zero if none exists.
\item[\textbf{Transcendental element.}] An element not algebraic over the base field, satisfying no nonzero polynomial with coefficients in the base field.
\item[\textbf{Irreducible polynomial.}] A nonconstant polynomial that cannot be factored as a product of two nonconstant polynomials over the base field.
\item[\textbf{Intermediate field.}] A field $E$ satisfying $K\subseteq E\subseteq L$; central to the Galois correspondence.
\end{description}
\section*{7. Sample Questions}
\emph{Exercise reference:} \cite{field_theory_ref1}. The cited edition is the source for the exercise set; consult its Exercises section for the official statement and numbering.
\begin{enumerate}[leftmargin=*,itemsep=0.9em]
\item \textbf{Exercise 1.} Find the degree $[\mathbb{Q}(\sqrt{2},\sqrt{3}):\mathbb{Q}]$ and a basis for $\mathbb{Q}(\sqrt{2},\sqrt{3})$ over $\mathbb{Q}$.

\emph{Solution.} Since $\sqrt{3}\notin\mathbb{Q}(\sqrt{2})$ (if $\sqrt{3}=a+b\sqrt{2}$ then $3=a^2+2b^2+2ab\sqrt{2}$, forcing $ab=0$ and yielding a contradiction), the extension $\mathbb{Q}(\sqrt{2},\sqrt{3})/\mathbb{Q}$ has degree $[\mathbb{Q}(\sqrt{2},\sqrt{3}):\mathbb{Q}(\sqrt{2})]\cdot[\mathbb{Q}(\sqrt{2}):\mathbb{Q}]=2\cdot 2=4$. A basis is $\{1,\sqrt{2},\sqrt{3},\sqrt{6}\}$.

\item \textbf{Exercise 2.} Find the minimal polynomial of $1+\sqrt{2}$ over $\mathbb{Q}$.

\emph{Solution.} Set $\alpha=1+\sqrt{2}$, so $\alpha-1=\sqrt{2}$ and $(\alpha-1)^2=2$, giving $\alpha^2-2\alpha-1=0$. The polynomial $x^2-2x-1$ is irreducible over $\mathbb{Q}$ (no rational roots by the rational root theorem), so it is the minimal polynomial.

\item \textbf{Exercise 3.} Find the multiplicative inverse of $3$ in $\mathbb{F}_7$ and of $x+1$ in $\mathbb{F}_2[x]/(x^3+x+1)$.

\emph{Solution.} In $\mathbb{F}_7$: $3\cdot 5=15\equiv 1\pmod{7}$, so $3^{-1}=5$. In $\mathbb{F}_2[x]/(x^3+x+1)$: perform polynomial long division of $x^3+x+1$ by $x+1$. We get $x^3+x+1=(x+1)(x^2+x)+1$, so $(x+1)(x^2+x)\equiv 1$, meaning $(x+1)^{-1}=x^2+x$.

\item \textbf{Exercise 4.} Let $f(x)=x^3-2$ over $\mathbb{Q}$. Prove that $[\mathbb{Q}(\sqrt[3]{2}):\mathbb{Q}]=3$.

\emph{Solution.} The polynomial $x^3-2$ is irreducible over $\mathbb{Q}$ by Eisenstein's criterion with $p=2$. Therefore its minimal polynomial over $\mathbb{Q}$ has degree $3$, so $[\mathbb{Q}(\sqrt[3]{2}):\mathbb{Q}]=3$.

\item \textbf{Exercise 5.} Prove that every finite field has order $p^n$ for some prime $p$ and positive integer $n$.

\emph{Solution.} A finite field $F$ has characteristic $p>0$ (since it is finite, $1+1+\cdots+1=0$ for some sum). Then $F$ contains $\mathbb{F}_p$ as a subfield. $F$ is a vector space over $\mathbb{F}_p$, say of dimension $n$. A basis $\{e_1,\ldots,e_n\}$ gives $|F|=p^n$, since every element is a unique $\mathbb{F}_p$-linear combination of the basis elements.

\end{enumerate}
\section*{8. References}
\addcontentsline{toc}{section}{References}

\begin{thebibliography}{9}
\bibitem{field_theory_ref1} David S. Dummit and Richard M. Foote, \emph{Abstract Algebra}, 3rd ed., Wiley, 2004.
\bibitem{field_theory_ref2} Ian Stewart, \emph{Galois Theory}, 4th ed., Chapman and Hall/CRC, 2015.
\bibitem{field_theory_ref3} Rudolf Lidl and Harald Niederreiter, \emph{Finite Fields}, 2nd ed., Cambridge University Press, 1997.
\bibitem{field_theory_ref4} Jean-Pierre Serre, \emph{Local Fields}, Springer, 1979.
\end{thebibliography}
